\section{Conclusion}\label{sec:conclusion}

In this paper, we presented SIREN (Stochastic It\=o Residual on Embedded Networks), a fundamentally novel framework for network intrusion detection. Moving away from the discrete-time, supervised classification paradigms that dominate current Graph Neural Network IDS research, SIREN conceptualizes the monitored network as a continuous-time dynamical system governed by graph-coupled It\=o Stochastic Differential Equations. By employing score-matching techniques to learn the continuous evolution of benign network traffic, our approach establishes a robust, unsupervised baseline representing normal operational states.

Crucially, SIREN introduces the \textit{Stein residual}—a continuous discrepancy metric comparing the model's predicted score with an empirical Kernel Density Estimate computed from live traffic. Aggregating these Stein residuals along the physical network topology amplifies the inherently weak anomaly signals generated by advanced persistent threats and coordinated lateral movement, allowing SIREN to detect complex exploits without requiring explicitly labelled attack datasets. 

Experimental evaluations on standard benchmarks, including UNSW-NB15 and CIC-IDS2018, demonstrate that SIREN effectively isolates systemic deviations from normal traffic distributions. Compared to state-of-the-art discrete-time GNN classifiers, our continuous-time approach exhibits superior robustness against adversarial perturbations and graph structural decay, while maintaining a strict, operationally viable False Alarm Rate. Future work will explore accelerating the Stein residual inference process using operator learning and extending the framework to handle highly heterogeneous IoT environments with sparse, bursty traffic profiles.
